% Reviewed statements of cubic_v3.tex, 2026-11-05. Not a standalone paper.
\begin{theorem}\label{thm:generalMonicCount} For every $\varepsilon>0$, as $H\to\infty$, the number of monic integral cubic polynomials with Galois group $C_3$ and root height at most $H$ is \[ H^3P(\log H)+O_\varepsilon(H^{2+\varepsilon}), \] where \begin{equation} P(X)=\frac{220\alpha}{243} \biggl(X+\mu-\frac{\zeta'(2)}{\zeta(2)}+ \beta +\log2-\frac{5}{88}\log3 +\frac{\pi\sqrt3}{10}-\frac{17}{15}\biggr) -\frac{20\pi\sqrt3}{81}. \end{equation} \end{theorem}

\begin{theorem}\label{thm:asymptoticPolynomialCount} For every $\varepsilon>0$, as $H\to\infty$, the number of monic integral trace-one cubic polynomials with Galois group $C_3$ and root height at most $H$ is \[ H^2P_1(\log H)+O_\varepsilon(H^{1+\varepsilon}) \] where \begin{equation} P_1(X)=\frac\alpha3 \biggl(X+\mu-\frac{\zeta'(2)}{\zeta(2)} + \beta +\frac78\log3-\frac12\biggr) -\frac{\pi}{3\sqrt3}. \end{equation} \end{theorem}

\begin{theorem}\label{thm:exactFormula} For $m\geq1$ let $a_E(m)$ denote the number of ideals of norm $m$ in the ring of integers $\mcO_E$ of $E$. If $t^3-t^2+at+b\in\ZZ[t]$ has Galois group $C_3$, then $a\leq0$. Fix $a \in \ZZ_{\leq 0}$ and put $n=1-3a$. The number of polynomials of the form $t^3 -t^2 + at + b \in \ZZ[t]$ with Galois group $C_3$ is equal to \begin{equation} \tfrac{1}{2}\big(a_E(n^3)-a_E(n)\big). \end{equation} \end{theorem}

\begin{equation}\label{eqn:splitForm} x_0x_1x_2=x_3^3. \end{equation}

\begin{equation}\label{eqn:surfaceQuadraticTwist} S^{(d)}\colon x_0(x_1^2-dx_2^2)=x_3^3. \end{equation}

\begin{theorem}\label{thm:heightZetaFunctionFormula} For $\mathrm{Re}(s)>2$ the series below converges absolutely, and \begin{equation} \sum_{P \in (\SSS\setminus\DDD_0)(\ZZ)} H(P)^{-s} = \left(1-\frac{1}{3^{s/2}}\right)^2 \xi\left(\frac s2\right)^2 \prod_{q\equiv2\pmod{3}}\left(1-\frac{1}{q^{s}}\right) \prod_{p\equiv1\pmod{3}}\left(1-\frac{3}{p^{s}}+\frac{2}{p^{3s/2}}\right) \end{equation} where $\xi$ is the Dedekind zeta function of $E$. This height zeta function can be meromorphically continued to the half-plane $\mathrm{Re}(s)>1$ and its only pole in this region is at $s = 2$ with order $2$. \end{theorem}

\begin{corollary}\label{cor:integralPointCount} For every $\varepsilon>0$, as $H\to\infty$, \[ \#\{P\in(\SSS\setminus\DDD_0)(\ZZ):H(P)\leq H\} =H^2Q(\log H)+O_\varepsilon(H^{1+\varepsilon}), \] where \begin{equation} Q(X)=\frac{2\alpha}{9} \biggl(X+\mu-\frac{\zeta'(2)}{\zeta(2)} +\beta+\frac38\log3-\frac12\biggr). \end{equation} \end{corollary}

\begin{equation}\label{eqn:polynomialWeights} w_f=\frac{|\mathrm{Aut}(C_3)|}{|\mathrm{Aut}(K_f,x_f)|}= \begin{cases} 1 &\text{if $f$ has a double root,}\\ 2 &\text{otherwise,} \end{cases} \end{equation}

\begin{theorem}[{\cite[\S2]{odesky2023moduli}}]\label{thm:modularInterpretation} The homogeneous variety $T$ is the moduli space for $C_3$-algebras with a given trace-one normal element. In particular, there is a bijection between rational points of $T$ and equivalence classes of $C_3$-algebras $K/\mathbb Q$ equipped with a trace one normal element $x \in K$. \end{theorem}

\begin{example}\label{example:twoRationalPointsSamePoly} Let $K$ be a cubic abelian number field. Then $K$, equipped with its canonical Galois action, is a $C_3$-algebra. The twist $K'$ of the $C_3$-algebra $K$ by the outer automorphism $g \mapsto g^{-1}$ of $C_3$ (with twisted action $g \ast x = g^{-1}x$) is not isomorphic to $K$ as a $C_3$-algebra.\footnote{ In terms of Galois cohomology, the non-cohomologous $1$-cocycles in $H^1(\QQ,C_3)$ corresponding to the $C_3$-algebras $K$ and $K'$ have the same image under the canonical map $H^1(\QQ,C_3) \to H^1(\QQ,S_3)$ because the outer automorphism of $C_3$ is realized by $S_3$-conjugation.} \end{example}

\begin{example}\label{example:splitAlgebra} Let $K_{\text{spl}} = \QQ^3$, the split cubic algebra. Then $C_3\subset S_3 = \mathrm{Aut}_{\text{$\QQ$-alg}}(K_{\text{spl}})$ and $K_{\text{spl}}$, equipped with its canonical $C_3$-action, is a $C_3$-algebra. Any transposition gives an isomorphism of $C_3$-algebras from $K_{\text{spl}}$ to its twist $K_{\text{spl}}'$. \end{example}

\begin{example}\label{example:splitAlgebraNormalElement} An element $x$ of trace one in the split $C_3$-algebra $K_{\text{spl}}$ is normal if and only if $x$ either has distinct coordinates or exactly two identical coordinates. For such a normal element $x$, the pairs $(K_{\text{spl}},x)$ and $(K_{\text{spl}}',x)$ are equivalent if and only if $x$ has exactly two identical coordinates (swapping the identical coordinates gives the required isomorphism); in particular, if $x$ has distinct coordinates then $(K_{\text{spl}},x)$ and $(K_{\text{spl}}',x)$ determine different rational points of $\mathcal G_1/C_3$, even though $K_{\text{spl}}$ and $K_{\text{spl}}'$ are isomorphic as $C_3$-algebras. \end{example}

\begin{proposition}\label{prop:basicfacts} Let $w$ be a place of $E$ lying over a place $v$ of $\QQ$. There is an exact sequence \begin{equation} 1 \longrightarrow K_v^{\mathrm{max}} \longrightarrow T(\QQ_v) \xlongrightarrow{\ell_w} \mathfrak t_w. \end{equation} If $w$ is infinite then $\ell_w$ is surjective, and if $w$ is finite then the image of $\ell_w$ is the lattice $X_\ast(T_{E})^{\Gamma(w/v)}$. \end{proposition}

\begin{equation}\label{eqn:heightforrationaladeles} H(x,\varphi) = \prod_{v \in Pl_\QQ} q_v^{-\frac{1}{e_v}\varphi(\ell_w(x_v)) } \end{equation}

\begin{lemma}\label{lemma:heightformula} Let $\mathbf 1$ denote the global section in $H^0(S_E,\mathcal O(3D_0))$ associated to the effective Cartier divisor $3D_0$, namely the image of $1$ under the natural inclusion $\mathcal O_{S_E}\hookrightarrow\mathcal O_{S_E}(3D_0)$. The space $H^0(S_E,\mathcal O(3D_0))$ is spanned over $E$ by the following four weight vectors: \begin{align}\label{eqn:weightVectors} &\mathbf 1,\qquad (1-3e_2e_1^{-2})\mathbf 1, \\ &(e_1^3-\tfrac 92 e_1 e_2+ \tfrac{27}{2} e_3 + \tfrac{\sqrt{-27}}{2} \sqrt{\mathrm{disc}}) e_1^{-3}\mathbf 1,\qquad (e_1^3-\tfrac 92 e_1 e_2+ \tfrac{27}{2} e_3 - \tfrac{\sqrt{-27}}{2} \sqrt{\mathrm{disc}}) e_1^{-3}\mathbf 1 \end{align} where $\sqrt{\mathrm{disc}} \coloneq (X-Z)(Y-X)(Z-Y)$. The associated characters of $T_E$ are, respectively, \begin{equation} 1,\quad \chi^{a'+b'},\quad \chi^{2a'+b'},\quad\chi^{2b'+a'} \end{equation} where $a' = 2v_1^\vee - v_2^\vee$ and $b' = 2v_2^\vee - v_1^\vee$ in the character lattice $M_E = X^\ast T_E$ and $(v_1^\vee,v_2^\vee)$ is the dual basis to $(v_1,v_2)$. \end{lemma}

\begin{equation}\label{eqn:formulaForNormCharacter} \chi^{a'+b'}(uv_0' + vv_1'+wv_2') = \frac{vw}{u^2} = \frac{(X+\zeta Y+\zeta^2 Z)(X+\zeta^2 Y+\zeta Z)}{(X+Y+Z)^2} = \frac{e_1^2-3e_2}{e_1^2} \end{equation}

\begin{equation}\label{eqn:formulaForNormCharacter2} \chi^{2b'+a'}(uv_0' + vv_1'+wv_2') =\frac{w^3}{u^3} = \frac{e_1^3-\tfrac 92 e_1 e_2+ \tfrac{27}{2} e_3 - \tfrac{\sqrt{-27}}{2} \sqrt{\mathrm{disc}}}{e_1^3} \end{equation}

\begin{lemma}\label{lemma:affineIntegrality} Write $e_2,e_3,\delta$ for the dehomogenized invariants \[ e_2=\frac{XY+YZ+ZX}{\varepsilon^2},\qquad e_3=\frac{XYZ}{\varepsilon^3},\qquad \delta=\frac{X^2Y+Y^2Z+Z^2X}{\varepsilon^3}. \] Then \begin{equation}\label{eqn:affineC3invariants} \ZZ[X/\varepsilon,Y/\varepsilon]^{C_3} =\ZZ[e_2,e_3,\delta]. \end{equation} A rational point $P$ of $S\setminus D_0$ is integral if and only if $e_2(P),e_3(P)\in\ZZ$. \end{lemma}

\begin{equation}\label{eqn:C3invariants} \ZZ[X,Y,Z]^{C_3} =\ZZ[e_1,e_2,e_3,X^2Y+Y^2Z+Z^2X] \end{equation}

\begin{proposition}\label{prop:canonicalHeight} The characteristic polynomial $f = t^3-t^2+at+b \in \QQ[t]$ of a rational point $(K/\mathbb Q,x) \in T(\mathbb Q)$ has integer coefficients if and only if $(K/\mathbb Q,x)$ is integral. For any integral rational point $(K/\QQ,x)$ on $T$, \begin{equation} H((K/\mathbb Q,x),\mathcal O(D_0)) = \sqrt{1-3a}. \end{equation} \end{proposition}

\begin{align}\label{eqn:localHeight} \max\Bigg( 1, |1-3e_2|_w, &\left| 1-\tfrac{9}{2} e_2+ \tfrac{27}{2} e_3 - \tfrac{\sqrt{-27}}{2} \sqrt{\mathrm{disc}} \right|_w, \\ &\left| 1-\tfrac{9}{2} e_2+ \tfrac{27}{2} e_3 + \tfrac{\sqrt{-27}}{2} \sqrt{\mathrm{disc}} \right|_w \Bigg)^{\frac{1}{[E:\QQ]}} \end{align}

\begin{align}\label{eqn:heightExtraFactor} \max\left( 1, |1-3e_2|^2, \left|1-\tfrac 92 e_2+ \tfrac{27}{2} e_3 + \tfrac{\sqrt{-27}}{2} \sqrt{\mathrm{disc}}\right|^2 \right)^{1/2}. \end{align}

\begin{lemma}[Root height and toric height]\label{lemma:rootVsToric} Let $P=(K/\QQ,x)$ be an integral rational point of $T$ with characteristic polynomial $f=t^3-t^2+at+b\in\ZZ[t]$. Write $H_{\mathrm{tor}}(f)=H(P,\mathcal O(D_0))$ for its toric height and set $h(f)=\max(|a|^{1/2},|b|^{1/3})$, so that $H(f)=\max\{1,h(f)\}$. Then $a\leq0$ and $h(f)^2=-a, H_{\mathrm{tor}}(f)^2=1+3h(f)^2$ so for every $B\geq0$, \begin{equation} h(f)\leq B \quad\Longleftrightarrow\quad H_{\mathrm{tor}}(f)\leq\sqrt{1+3B^2}. \end{equation} In particular, $H_{\mathrm{tor}}(f)\asymp H(f)$. The only polynomial here with $h(f)=0$ is $t^3-t^2$, for which $H_{\mathrm{tor}}(f)=1$. \end{lemma}

\begin{equation}\label{eqn:isoCocharLattices} N_E=\ZZ\langle\omega,\gamma\omega \rangle \to N_E'=\ZZ\langle v_1,v_2\rangle \end{equation}

\begin{proposition}[Parametrization of monic abelian cubic trace-one polynomials by $\QQ(\sqrt{-3})^\times$]\label{prop:cyclotomicParametrization} The tori $T$ and $\mathcal G_1 = R^E_\QQ \mathcal \GG_m$ are isomorphic as algebraic groups over $\QQ$. Every rational point $(K/\QQ,x)$ of $T$ thereby determines an element of $\QQ(\sqrt{-3})^\times$ which is canonically determined up to the action of $\mathrm{Aut}(\mathcal G_1)$. For integral rational points of $T$, the toric height $H_{\mathrm{tor}}(f)\coloneqq\sqrt{1-3a}$ is identified with the square-root of the norm on $\QQ(\sqrt{-3})^\times$. Let $\zeta$ be a primitive cube root of unity. If $u+v\zeta \in \QQ(\sqrt{-3})^\times$ has norm $N$ and trace $T$, then the characteristic polynomial of the corresponding rational point $(K/\QQ,x)$ is \begin{equation} f=t^3-t^2+\tfrac{1}{3}(1 - N)t+\tfrac{1}{27}(N(3-T)-1) \in \QQ[t]. \end{equation} Such a polynomial either has Galois group $C_3$ or splits into three linear factors over $\QQ$, with at most two linear factors being the same. Conversely, a monic trace-one polynomial $f=t^3-t^2+at+b \in \QQ[t]$ which either has Galois group $C_3$ or splits into three linear factors over $\QQ$, with at most two linear factors being the same, can be expressed in this way for precisely two rational points of $T$ if $f$ has no repeated roots, or for precisely one rational point of $T$ if $f$ has a double root which is not a triple root. The elements $u+v\zeta \in \QQ(\sqrt{-3})^\times$ corresponding to $f=t^3-t^2+at+b$ will be the roots of the quadratic polynomial \begin{equation} g = t^2 - \left(3-\frac{1+27b}{1-3a}\right)t + 1-3a \in \QQ[t]. \end{equation} The polynomial $f$ will have integral coefficients if and only if \begin{equation}\label{eqn:integralityConditionsUsingCyclo} \begin{cases} u^2+v^2-uv \in 1+3 \ZZ \text{ and }\\ (u^2+v^2-uv)(3-2u+v) \in 1+27 \ZZ. \end{cases} \end{equation} \end{proposition}

\begin{theorem}\label{thm:heightZetaFunctionFormula2} Fix any $s \in (\CC^{\Sigma(1)})^{\Gamma}$ with $\mathrm{Re}(s_0+s_1+s_2) > 2$. Then the multivariate Dirichlet series \begin{equation} Z(s)= \sum_{\substack{P \in T(\mathbb Q)\\\text{\emph{integral}}}} H(P,s)^{-1} \end{equation} is absolutely convergent and equals \begin{align} \left({1-3^{-z}}\right) \zeta(z) \prod_{q\equiv2\pmod{3}} \bparen{1+\frac{1}{q^{z}}}^{-1} \prod_{p\equiv1\pmod{3}} \bparen{1+\frac{3}{p^z}\bparen{1-\frac{1}{p^z}}^{-1}} \end{align} where $z = \frac{1}{2}(s_0+s_1+s_2)$. This multivariate Dirichlet series admits a meromorphic continuation to $\{s \in (\CC^{\Sigma(1)})^{\Gamma} : \mathrm{Re}(s_0+s_1+s_2) > 1\}$. \end{theorem}

\begin{proposition}\label{prop:localIntegrality} Let $v$ be a finite place of $\QQ$ and let $w$ be a place of $E$ lying over $v$. \begin{enumerate} \item If $v \neq 3$, then $1_v^{D_0}(x) = 1$ if and only if $\ell_w(x) \in \mathbb R_{\geq 0}v_1+\mathbb R_{\geq 0}v_2$. In particular $K_v^{D_0}=K_v^{\mathrm{max}}$. \item Let $v = 3$, let $w$ be the unique place of $E$ over $3$, and let $\mathfrak w = (1-\zeta)O_{E,w}$ be the maximal ideal of $O_{E,w}$, so that $3O_{E,w} = \mathfrak w^2$. The support of $1_3^{D_0}$ is the subgroup \begin{equation} K_{3,2}\coloneqq\iota^{-1}(1+\mathfrak w^2) =\iota^{-1}(1+3O_{E,w})\subset K_3^{\mathrm{max}}. \end{equation} \item Let $K^{\mathrm{max}}=\prod_v K_v^{\mathrm{max}}$ and $K^{D_0}=\prod_v K_v^{D_0}\subset K^{\mathrm{max}}$. The latter is the subgroup which fixes the characteristic function $1^{D_0}=\otimes_v1_v^{D_0}$ for integral points in $T(\AA)$, with $1_\infty^{D_0}=1$. Then \[ K_3^{D_0}=K_{3,2},\qquad K^{D_0}=K_{3,2}\times\prod_{v\neq3}K_v^{\mathrm{max}}. \] \end{enumerate} \end{proposition}

\begin{equation}\label{eqn:poissonae} \int_{T(\QQ)} f(xy)\, dx = \int_{T(\QQ)^\perp} \widehat{f}(\chi) \chi(y) \,d\chi \end{equation}

\begin{proposition}\label{prop:heightFourierTransform} Let $s \in (\CC^{\Sigma(1)})^{\Gamma}$ and assume $\mathrm{Re}(s_e)>0$ for each $e \in\{0,1,2\}$. Let $w$ be the infinite place of $E$. Let $\chi\in T(\RR)^\vee$ be a unitary character. If $\chi$ is ramified then $\widehat{H_\infty}(\chi,-s,D_0)$ is identically zero. If $\chi$ is unramified, then $\chi(x)= e(\langle \ell_w(x) ,m\rangle)$ for all $x \in T(\RR)$ for a unique $m \in \mathfrak t^\vee$, and \begin{equation}\label{eqn:infiniteHeightFourierTransform} \widehat{H_\infty}(\chi,-s,D_0) = \left(\frac{-1}{2\pi i}\right) \frac{s_0+s_1+s_2}{2\pi i} \frac{1}{(m(v_0)+\frac{s_0}{2\pi i})(m(v_0)-\frac{s_1+s_2}{2\pi i})}. \end{equation} Next let $v$ be a finite place of $\QQ$. For any unitary character $\chi\in T(\QQ_v)^\vee$, the integral defining $\widehat{H_v}(\chi,-s,D_0)$ converges absolutely to a holomorphic function of $s$ in the region $$\{s \in (\CC^{\Sigma(1)})^\Gamma : \mathrm{Re}(s_1),\mathrm{Re}(s_2)>0 \}.$$ Assume $v \neq 3$. Let $w$ be any place of $E$ lying over $v$. The local characteristic function $1_v^{D_0}$ is $K_v^{\mathrm{max}}$-invariant. If $\chi$ is ramified, then $\widehat{H_v}(\chi,-s,D_0)$ is identically zero. If $\chi$ is unramified then we may regard $\chi$ as a character on $X_\ast(T_{E})^{\Gamma(w/v)}$ (Proposition~\ref{prop:basicfacts}) and \begin{equation}\label{eqn:finiteHeightFourierTransform} \widehat{H_v}(\chi,-s,D_0) = \sum_{\substack{n \in X_\ast(T_{E})^{\Gamma(w/v)}\\ n \in \mathbb R_{\geq 0}v_1+\mathbb R_{\geq 0}v_2} } \chi(n)^{-1} q_v^{\varphi(n)} . \end{equation} If $v = 3$ then the support of $x \mapsto H_3(x,-s,D_0)$ is the congruence subgroup $K_{3,2}$ of $K_3^{\mathrm{max}}$ of index six. Under the isomorphism $\iota:T(\QQ_3)\to E_w^\times$ the support corresponds to the subgroup $1+3O_{E,w}$ of $O_{E,w}^\times$ where $w$ is the unique place of $E$ lying over $3$. \end{proposition}

\begin{equation}\label{eqn:localfourierassumovern} \int_{T(\QQ_v)} H_v(x,s)^{-1}1_v^{D_0}(x)\chi(x)^{-1}\,d^\times x_v =\sum_{n \in N_w} q_v^{\frac{1}{e_v}\varphi(n)} 1_v^{D_0}(n) \chi(n)^{-1} = \sum_{\substack{n \in X_\ast(T_{E})^{\Gamma(w/v)}\\ n \in \mathbb R_{\geq 0}v_1+\mathbb R_{\geq 0}v_2} } \chi(n)^{-1} q_v^{\varphi(n)} . \end{equation}

\begin{equation}\label{eqn:LAsNorm} \ell(x)(\mathrm{N}) = -\log |\mathrm{N}(x)|_{\AA} \end{equation}

\begin{proposition}\label{prop:Lmap} There is an exact sequence \begin{equation} 1 \longrightarrow K^{\mathrm{max}}/\mu^{\mathrm{max}} \longrightarrow T(\mathbb Q)\backslash T(\AA) \xlongrightarrow{\ell} \mathfrak t \longrightarrow 0 \end{equation} where $K^{\mathrm{max}}$ is the maximal compact subgroup of $T(\AA)$ and $\mu^{\mathrm{max}} = T(\QQ) \cap K^{\mathrm{max}}$. \end{proposition}

\begin{lemma}\label{lemma:KprimeGroup} Let $K^{D_0}=K_{3,2}\times\prod_{v\neq3}K_v^{\mathrm{max}}$ be the subgroup of $K^{\mathrm{max}}$ fixing $1^{D_0}$, as in Proposition~\ref{prop:localIntegrality}. The map $\ell$ induces an isomorphism \begin{equation} T(\QQ)\backslash T(\AA)/K^{D_0} \xlongrightarrow{\sim} \mathfrak t. \end{equation} Let $a(u)\in T(\AA)$ have infinite component $\iota^{-1}(e^u)$ and all finite components one, so $\ell(a(u))=uv_0$. In this coordinate, the pushforward of the quotient measure is $du/6$. \end{lemma}

\begin{equation}\label{eqn:chimdefn} \chi_m(x)=e(\langle m,\ell(x)\rangle), \qquad m\in\mathfrak t^\vee. \end{equation}

\begin{lemma}[{Finite Fourier product}]\label{lemma:finiteFourierProduct} For $\operatorname{Re}(\sigma)>1$ and $t\in\RR$, put $\lambda=\sigma-\pi it$. Then \begin{equation}\label{eqn:finiteSpectralProduct} \widehat{H_{\mathrm{fin}}}(\chi_m,-s,D_0) =\frac16 \prod_{q\equiv2\pmod3}(1-q^{-2\lambda})^{-1} \prod_{p\equiv1\pmod3} \frac{1+2p^{-\lambda}}{(1-p^{-\lambda})^2} \end{equation} and \begin{equation}\label{eqn:finiteEulerBound} \left|\widehat{H_{\mathrm{fin}}}(\chi_m,-s,D_0)\right| \leq\widehat{H_{\mathrm{fin}}}(\chi_0,{-\operatorname{Re}(s)},D_0)<\infty. \end{equation} Here $\chi_0$ is the trivial character, corresponding to $m=0$. The Euler product converges absolutely and locally uniformly for $\operatorname{Re}(\lambda)>1$ and is bounded on every half-plane $\operatorname{Re}(\lambda)\geq\tau>1$. \end{lemma}

\begin{lemma}\label{lemma:heightZetaConvergence} Let $s=(s_0,\sigma,\sigma)\in(\CC^{\Sigma(1)})^\Gamma$ with $\operatorname{Re}(s_0)>0$ and $\operatorname{Re}(\sigma)>1$, and put $f=H(\cdot,-s,D_0)$. Then $f\in L^1(T(\AA))$ and $\widehat f\,|_{T(\QQ)^\perp}\in L^1(T(\QQ)^\perp)$. \end{lemma}

\begin{equation}\label{eqn:archimedeanSpectralFactor} \begin{aligned} \widehat{H_\infty}(\chi_{m,\infty},-s,D_0) &=\frac1{s_0+2\pi it}+\frac1{2\sigma-2\pi it}\\ &=\frac{s_0+2\sigma} {(s_0+2\pi it)(2\sigma-2\pi it)}. \end{aligned} \end{equation}

\begin{lemma}\label{lemma:poissonIntegral} For $s$ as in Lemma~\ref{lemma:heightZetaConvergence}, we have \begin{equation}\label{eqn:poissonEulerIntegral} Z(s)=6\int_\RR \widehat{H_\infty}(\chi_{m,\infty},-s,D_0) \widehat{H_{\mathrm{fin}}}(\chi_m,-s,D_0)\,dt, \end{equation} where both the series defining $Z(s)$ and the integral converge absolutely. \end{lemma}

\begin{equation}\label{eqn:normCharacterDivisor} \operatorname{div}(\chi^{a'+b'}) =\sum_{i=0}^2\langle a'+b',v_i\rangle D_i =-2D_0+D_1+D_2. \end{equation}

\begin{equation}\label{eqn:eulerCorrectionFactor} R(z)=\prod_{q\equiv2\pmod3}\left(1-q^{-2z}\right) \prod_{p\equiv1\pmod3}\left(1-3p^{-2z}+2p^{-3z}\right). \end{equation}

\begin{lemma}\label{lemma:simplifiedZ(s)} The height zeta function is also given by \begin{align} Z(s)=\left(1-3^{-z}\right)^2\xi(z)^2R(z),\qquad z=\tfrac12(s_0+s_1+s_2), \end{align} where $\xi$ is the Dedekind zeta function of the cyclotomic field $\QQ(\sqrt{-3})$. The height zeta function has meromorphic continuation to the region $\{s \in (\CC^{\Sigma(1)})^{\Gamma} : \mathrm{Re}(s_0+s_1+s_2) > 1\}$. \end{lemma}

\begin{equation}\label{eqn:Z0factorization} Z_0(z)=\left(1-3^{-z}\right)^2\xi(z)^2R(z). \end{equation}

\begin{proposition}\label{prop:leadingTerms} The only pole of $Z_0(z)$ in the half-plane $\mathrm{Re}(z)>1/2$ is a double pole at $z=1$. Its Laurent expansion there is \begin{equation}\label{eqn:Z0Laurent} Z_0(z)=\frac{\alpha}{9(z-1)^2} +\frac{2\alpha}{9}\left(\mu-\frac{\zeta'(2)}{\zeta(2)}+\beta +\frac38\log3\right)\frac{1}{z-1}+O(1), \end{equation} where $\alpha$ and $\beta$ are as in \eqref{eqn:introConstants}. \end{proposition}

\begin{lemma}\label{lemma:reducibles} Let $d_n$ be the coefficients of $Z_0(z)=\sum_{n\geq1}d_n n^{-z}$. For $X\geq1$, \begin{equation}\label{eqn:weightedPolynomialCount} \sum_{\substack{f\in F\\H_{\mathrm{tor}}(f)\leq\sqrt X}}1 =\frac12\sum_{n\leq X}d_n -\frac{\pi}{9\sqrt3}X+O(\sqrt X). \end{equation} \end{lemma}

\begin{equation}\label{eqn:heightMeanBound} \int_U^{2U}|Z_0(\sigma+it)|\,dt \leq A_\eta\int_U^{2U}|\xi(\sigma+it)|^2\,dt \qquad (U \geq 2). \end{equation}

\begin{equation}\label{eqn:irrPolys} \sum_{\substack{f\in F\\H_{\mathrm{tor}}(f)\leq\sqrt X}}1 =\frac12XQ\left(\tfrac12\log X\right) -\frac{\pi}{9\sqrt3}X +O_\varepsilon(X^{1/2+\varepsilon}). \end{equation}

\begin{lemma}\label{lemma:generalHeightReduction} For $s\in\ZZ$, define $\phi_s$ on $F_s$ by \[ \phi_s(f)(t)= \begin{cases} f(t+\tfrac{s}{3}),&s\equiv0\pmod3,\\ f(t+\tfrac{s-1}{3}),&s\equiv1\pmod3,\\ -f(-t+\tfrac{s+1}{3}),&s\equiv-1\pmod3. \end{cases} \] Then $\phi_s$ is a bijection from $F_s$ to $F_0$ if $3\mid s$, and from $F_s$ to $F_1$ if $3\nmid s$, and it preserves $n$. For $f\in F_s$, put $n=n(f)$. Then $n>0$, and \[ H(f)\leq H \quad\Longleftrightarrow\quad |s|\leq H\ \text{ and }\ n\leq s^2+3H^2. \] If $3\nmid s$, then $H_{\mathrm{tor}}(\phi_s(f))^2=n$. If $3\mid s$, then $\phi_s(f)=t^3-\frac n3t+c$ for some $c\in\ZZ$. \end{lemma}

\begin{equation}\label{eqn:A1asymptotic} A_1(X) =\frac12XQ\left(\tfrac12\log X\right) -\frac{\pi}{9\sqrt3}X +O_\varepsilon(X^{1/2+\varepsilon}). \end{equation}

\begin{equation}\label{eqn:dSummatoryForGeneralTrace} \sum_{n\leq X}d_n =XQ\left(\tfrac12\log X\right) +O_\varepsilon(X^{1/2+\varepsilon}), \end{equation}

\begin{lemma}\label{lemma:traceZeroCoefficients} Let $e$ be the exponent of $3$ in $N$, and write $N=3^eM$ with $3\nmid M$. Then \[ r(N)= \begin{cases} d_M,&e=0,\\ 3d_M+1,&N=3u^2\text{ for some }u\in\ZZ_{>0},\\ 3d_M,&\text{otherwise}. \end{cases} \] \end{lemma}

\begin{proposition}\label{prop:traceZeroAsymptotic} For every $\varepsilon>0$, \[ A_0(X) =\frac52XQ\left(\tfrac12\log X\right) -\left(\frac{\alpha}{4}\log3+\frac{\pi}{3\sqrt3}\right)X +O_\varepsilon(X^{1/2+\varepsilon}). \] \end{proposition}

\begin{equation}\label{eqn:monicArchimedeanMoments} \begin{aligned} \int_{-1}^1(3+u^2)\,du&=\frac{20}{3},\\ \int_{-1}^1(3+u^2)\log(3+u^2)\,du &=\frac{40}{3}\log2+\frac{4\pi\sqrt3}{3}-\frac{76}{9}. \end{aligned} \end{equation}

\begin{align} \#\{f\in F_{\mathrm{mon}}:H(f)\leq H\} =\sum_{\substack{s\in\ZZ,\ |s|\leq H\\3\nmid s}} A_1(3H^2+s^2) +\sum_{\substack{s\in\ZZ,\ |s|\leq H\\3\mid s}} A_0\left(H^2+\frac{s^2}{3}\right). \label{eqn:generalTraceDecomposition} \end{align}

\begin{rem}\label{rem:translationClasses} We derive the count of translation classes by $H_{\mathrm{BS}}$ stated in the introduction. The invariant $I$ there is $n(f)$. Translating the roots by an integer $k$ changes the trace by $3k$ and preserves $n(f)$, so each translation class in $F_{\mathrm{mon}}$ has exactly one representative with trace in $\{-1,0,1\}$. The map $f(t)\mapsto-f(-t)$ negates the roots, so it matches the classes of trace $1$ with those of trace $-1$ and preserves $n(f)$. A trace-one polynomial $t^3-t^2+at+b$ has $n(f)=1-3a$, and $t^3-Nt+c$ has $n(f)=3N$. Hence the number $C(Y)$ of translation classes with $I\leq Y$ is \[ C(Y)=2A_1(Y)+A_0(Y/3). \] Substituting \eqref{eqn:A1asymptotic} and Proposition~\ref{prop:traceZeroAsymptotic}, and using that $Q$ is linear with leading coefficient $2\alpha/9$, gives \[ C(Y)=Y\left(\frac{11}{6}Q\left(\tfrac12\log Y\right) -\frac{19\alpha}{108}\log3-\frac{\pi}{3\sqrt3}\right) +O_\varepsilon(Y^{1/2+\varepsilon}). \] Since $f$ has three distinct real roots, $4I^3-J^2=27\operatorname{disc}(f)>0$, so $H_{\mathrm{BS}}(f)=I^3$. Setting $Y=H^{1/3}$ and substituting the formula for $Q$ from Corollary~\ref{cor:integralPointCount} gives the stated polynomial $P_{\mathrm{BS}}$ and error $O_\varepsilon(H^{1/6+\varepsilon})$. \end{rem}

\begin{theorem}\label{thm:traceDistribution} Uniformly in $r$, $u_1$ and $u_2$, as $H\to\infty$, \[ N_r(H;u_1,u_2) =c_r\left(\int_{u_1}^{u_2}(3+u^2)\,du\right)H^3\log H+O(H^3), \qquad c_0=\frac{5\alpha}{81},\quad c_1=c_2=\frac{\alpha}{27}. \] Consequently, among the polynomials counted in Theorem~\ref{thm:generalMonicCount}, the proportions with trace congruent to $0$, $1$ and $2$ modulo $3$ tend to $5/11$, $3/11$ and $3/11$. Within each residue class, $\operatorname{Tr}(f)/H$ has limiting density $\frac{3}{20}(3+u^2)$ on $[-1,1]$. \end{theorem}

\begin{equation}\label{eqn:monicConstantsFromTraceCounts} P(X) =\frac{110}{27}Q(X) +\frac{220\alpha}{243}\left(\log2-\frac{19}{44}\log3 -\frac{19}{30}+\frac{\pi\sqrt3}{10}\right) -\frac{20\pi\sqrt3}{81}. \end{equation}

\begin{proposition}[Integral cubic model]\label{prop:surfaceIntegralCubic} The invariant cubics \[ A=e_1^3,\qquad B=e_1e_2,\qquad C=e_3,\qquad D=\delta \] identify $\SSS$ with the cubic hypersurface in $\PP^3_\ZZ$ defined by \begin{equation}\label{eqn:surfaceIntegralCubic} F(A,B,C,D):=B^3+AD^2-ABD+3ACD+A^2C-6ABC+9AC^2=0. \end{equation} The divisor $\DDD_0=\{e_1=0\}$ corresponds to the line $\{A=B=0\}$ in these coordinates. On the generic fiber, \begin{equation}\label{eqn:surfaceAnticanonical} 3D_0\sim-K,\qquad \mcO(3D_0)\cong\mcO(1),\qquad (-K)^2=3. \end{equation} In particular, this cubic embedding is anticanonical. \end{proposition}

\begin{equation}\label{eqn:surfaceInvariantRelation} \delta^2-(e_1e_2-3e_3)\delta +e_2^3+e_1^3e_3-6e_1e_2e_3+9e_3^2=0. \end{equation}

\begin{proposition}\label{prop:surfaceRationalForm} Over $\ZZ[1/6]$, $\SSS$ is isomorphic to the cubic hypersurface \begin{equation}\label{eqn:surfaceNormCubic} x_0(x_1^2+3x_2^2)=x_3^3 \end{equation} in $\PP_{\ZZ[1/6]}^3$. Its generic fiber $S$ is a del Pezzo surface of degree $3$ and geometric Picard rank $1$. Its three geometric singularities are of type $\mathbf A_2$: one is rational, and the other two are conjugate over $E=\QQ(\sqrt{-3})$. Over $E$, the coordinates $x_1'=x_1+\sqrt{-3}x_2$ and $x_2'=x_1-\sqrt{-3}x_2$ give the split toric cubic $x_0x_1'x_2'=x_3^3$. In these coordinates, $D_0=\{x_0=x_3=0\}$. \end{proposition}

\begin{equation}\label{eqn:surfaceCoordinateChange} \begin{aligned} x_0&=A,& x_3&=A-3B,\\ x_1&=A-\tfrac92B+\tfrac{27}{2}C,& x_2&=3D-\tfrac32B+\tfrac92C. \end{aligned} \end{equation}

\begin{proposition}[Resolution and reduced boundary]\label{prop:surfaceResolvedBoundary} The surface $S$ admits a minimal resolution $\pi\colon X\to S$ over $\QQ$. Put $D_{\mathrm{red}}=(\pi^{-1}D_0)_{\mathrm{red}}$. Then $X$ is a weak del Pezzo surface of degree $3$, of geometric Picard rank $7$. The geometric components of $D_{\mathrm{red}}$ form a chain with self-intersections \begin{equation}\label{eqn:surfaceBoundaryChain} (-2)-(-2)-(-1)-(-2)-(-2). \end{equation} The divisor $L=-K_X-D_{\mathrm{red}}$ is nef and not numerically zero, but is not big. More precisely, \begin{equation}\label{eqn:surfaceLogIntersections} K_X\cdot D_{\mathrm{red}}=-1,\qquad D_{\mathrm{red}}^2=-1,\qquad L^2=0. \end{equation} \end{proposition}

\begin{proposition}\label{prop:santensConstant} Let $\pi\colon X\to S$ be the minimal resolution, let $D\subset X$ be the reduced preimage of $D_0$, and put $U=X\setminus D$. The coordinate ring of the affine chart $\SSS\setminus\DDD_0$ pulls back to a finitely generated $\ZZ$-subalgebra \[ A:=\pi^*\Gamma(\SSS\setminus\DDD_0,\mathcal O) \subset\Gamma(U,\mathcal O). \] A point $P\in U(\QQ)$ is integral if $g(P)\in\ZZ$ for every $g\in A$. With this integral condition and the height $H(\pi(-),1)$, \cite[Theorem~8.7]{santens2023integral} gives \[ \#\{P\in(\SSS\setminus\DDD_0)(\ZZ):H(P,1)\leq B\} \sim cB^2\log B,\qquad c=\frac{2\alpha}{9}. \] \end{proposition}

\begin{equation}\label{eqn:santensLocalHeight} H_v(P,3)=\prod_{w\mid v} \max\bigl(1,|n|_w,|\nu^2\overline\nu|_w, |\nu\overline\nu^2|_w\bigr)^{1/2}. \end{equation}

\begin{equation}\label{eqn:santensTau} c=\frac{\theta}{a(b-1)!}\,\tau=\tau,\qquad \tau=2\rho_E\prod_v\tau'_v, \end{equation}

\begin{equation}\label{eqn:introConstants} \mu=\lim_{s\to1} \left(\frac{\xi'(s)}{\xi(s)}+\frac{1}{s-1}\right), \quad \alpha=\prod_{p\equiv1\pmod3}\left(1-\frac{2}{p(p+1)}\right), \quad \beta=\sum_{p\equiv1\pmod3}\frac{(2p+1)\log p}{(p^2-1)(p+2)}. \end{equation}

\begin{equation}\label{eqn:Z0LaurentR} Z_0(z)=\frac{a}{(z-1)^2} +2a\left(\mu+\frac12\log3+\frac{R'(1)}{2R(1)}\right)\frac{1}{z-1} +O(1), \qquad a=\frac49r^2R(1). \end{equation}

\begin{lemma}\label{lemma:normalityCriterion} Let $K/\QQ$ be a $C_3$-algebra and let $x\in K$ have characteristic polynomial $t^3-e_1t^2+e_2t-e_3$. Then $x$ is normal if and only if $e_1(e_1^2-3e_2)\neq0$. \end{lemma}

\begin{rem}\label{rem:heightCoefficients} Comparing Euler factors in \eqref{eqn:Z0factorization} gives the coefficients explicitly: $d_n=0$ if $3\mid n$, and if $3\nmid n$ then \[ d_n=\sum_{d\mid n}3^{\omega(d_1)}(-1)^{\Omega(d_2)}, \] where $d=d_1d_2$ with $d_1$ divisible only by primes $p\equiv1\pmod3$ and $d_2$ only by primes $q\equiv2\pmod3$. Equivalently, $d_n$ is the product of $3k+1$ over prime powers $p^k\parallel n$ with $p\equiv1\pmod3$ if every prime $q\equiv2\pmod3$ divides $n$ to an even power, and $d_n=0$ otherwise. \end{rem}

\begin{lemma}\label{lemma:xiMeanValue} Let $\eta>0$. For $1/2\leq\sigma\leq2$ and $U\geq2$, \begin{equation}\label{eqn:xiMeanBound} \int_U^{2U}|\xi(\sigma+it)|^2\,dt \ll_\eta U^{1+\eta}. \end{equation} \end{lemma}

\begin{equation}\label{eqn:criticalMeanSquare} \int_1^T|\xi(1/2+it)|^2\,dt \ll T\log^2T. \end{equation}

\begin{rem}\label{rem:localNonintegrability} The local Fourier transforms --- and therefore the entire Poisson summation argument --- must be computed \emph{before} restricting to the line in $\mathrm{Pic}^T(S) \otimes \CC$ spanned by the $T_E$-line bundle $\mathcal O(D_0)$ of interest since $x \mapsto H_v(x,-s,D_0)$ will not be integrable for any place $v \neq 3,\infty$ once either of $s_1$ or $s_2$ vanishes, no matter how large and positive $\mathrm{Re}(s_0)$ is. \end{rem}

\begin{rem}\label{rem:idealParametrization} We prompted GPT-$6.1$ with the contents of an earlier version of this paper and asked it to identify the Dirichlet series in Theorem~\ref{thm:heightZetaFunctionFormula}. Using our Eisenstein parametrization (\S\ref{sec:quadraticParametrization}) it proved that \[ \sum_{P \in (\SSS\setminus\DDD_0)(\ZZ)} H(P)^{-s} = \sum_{\substack{0\ne J\subseteq\mcO_E\\ J+3\mcO_E=\mcO_E\\ \mathrm N J\text{ a cube}}} N^E_\QQ(J)^{-s/6} \] where $J \subset \mcO_E$ runs over integral nonzero ideals. It also found a result in the literature \cite{kuehleitnerNowak2004} that we previously missed. After translating the $C_3$-polynomial counting problem into an integral Diophantine problem on $S^{(-3)}$ (explained in the next section), K\"uhleitner--Nowak's~\cite[Corollary~2]{kuehleitnerNowak2004} can be used to obtain an shorter proof of Corollary~\ref{cor:integralPointCount} and Theorem~\ref{thm:asymptoticPolynomialCount} without using the toric method, and with an error term $O\bigl(H(\log H)^3(\log\log H)^2\bigr)$. We did not incorporate this argument as the toric approach has the merit of generalizing readily to other abelian groups where there is no analogue of \cite{kuehleitnerNowak2004}. \end{rem}
